\documentclass[10pt,a4paper]{amsart}
\usepackage{amsmath,amssymb,amsthm,graphicx}
% For dealing with references we use the comment environment
\usepackage{verbatim}
\newenvironment{reference}{\comment}{\endcomment}
%\newenvironment{reference}{}{}
\newenvironment{slogan}{\comment}{\endcomment}
\newenvironment{history}{\comment}{\endcomment}
% For commutative diagrams we use Xy-pic
\usepackage[all]{xy}

% We use 2cell for 2-commutative diagrams.
\xyoption{2cell}
\UseAllTwocells
\theoremstyle{plain}
\newtheorem{theorem}[subsection]{Theorem}
\newtheorem{proposition}[subsection]{Proposition}
\newtheorem{lemma}[subsection]{Lemma}

\theoremstyle{definition}
\newtheorem{definition}[subsection]{Definition}
\newtheorem{example}[subsection]{Example}
\newtheorem{exercise}[subsection]{Exercise}
\newtheorem{situation}[subsection]{Situation}

\theoremstyle{remark}
\newtheorem{remark}[subsection]{Remark}
\newtheorem{remarks}[subsection]{Remarks}

\numberwithin{equation}{subsection}
% Macros
%
\def\lim{\mathop{\mathrm{lim}}\nolimits}
\def\colim{\mathop{\mathrm{colim}}\nolimits}
\def\Spec{\mathop{\mathrm{Spec}}}
\def\Hom{\mathop{\mathrm{Hom}}\nolimits}
\def\Ext{\mathop{\mathrm{Ext}}\nolimits}
\def\SheafHom{\mathop{\mathcal{H}\!\mathit{om}}\nolimits}
\def\SheafExt{\mathop{\mathcal{E}\!\mathit{xt}}\nolimits}
\def\Sch{\mathit{Sch}}
\def\Mor{\mathop{\mathrm{Mor}}\nolimits}
\def\Ob{\mathop{\mathrm{Ob}}\nolimits}
\def\Sh{\mathop{\mathit{Sh}}\nolimits}
\def\NL{\mathop{N\!L}\nolimits}
\def\CH{\mathop{\mathrm{CH}}\nolimits}
\def\proetale{{pro\text{-}\acute{e}tale}}
\def\etale{{\acute{e}tale}}
\def\QCoh{\mathit{QCoh}}
\def\Ker{\mathop{\mathrm{Ker}}}
\def\Im{\mathop{\mathrm{Im}}}
\def\Coker{\mathop{\mathrm{Coker}}}
\def\Coim{\mathop{\mathrm{Coim}}}

% Boxtimes
%
\DeclareMathSymbol{\boxtimes}{\mathbin}{AMSa}{"02}

%
% Macros for moduli stacks/spaces
%
\def\QCohstack{\mathcal{QC}\!\mathit{oh}}
\def\Cohstack{\mathcal{C}\!\mathit{oh}}
\def\Spacesstack{\mathcal{S}\!\mathit{paces}}
\def\Quotfunctor{\mathrm{Quot}}
\def\Hilbfunctor{\mathrm{Hilb}}
\def\Curvesstack{\mathcal{C}\!\mathit{urves}}
\def\Polarizedstack{\mathcal{P}\!\mathit{olarized}}
\def\Complexesstack{\mathcal{C}\!\mathit{omplexes}}
% \Pic is the operator that assigns to X its picard group, usage \Pic(X)
% \Picardstack_{X/B} denotes the Picard stack of X over B
% \Picardfunctor_{X/B} denotes the Picard functor of X over B
\def\Pic{\mathop{\mathrm{Pic}}\nolimits}
\def\Picardstack{\mathcal{P}\!\mathit{ic}}
\def\Picardfunctor{\mathrm{Pic}}
\def\Deformationcategory{\mathcal{D}\!\mathit{ef}}
\usepackage{cleveref}
\crefname{originalreference}{original reference}{original references}
\Crefname{originalreference}{Original reference}{Original references}
\makeatletter
\newlabel{situation-algebraize}{{0EHC}{source}{}{sourcepointer-1}{}}
\newlabel{situation-algebraize@cref}{{[originalreference][0][]0EHC}{[1][0][]source}{}{}{}}
\makeatother
\crefname{originaldefinition}{original definition}{original definitions}
\Crefname{originaldefinition}{Original definition}{Original definitions}
\makeatletter
\newlabel{definition-s-d-inequalities}{{0EJ3}{source}{}{sourcepointer-2}{}}
\newlabel{definition-s-d-inequalities@cref}{{[originaldefinition][0][]0EJ3}{[1][0][]source}{}{}{}}
\makeatother
\crefname{originalequation}{original equation}{original equations}
\Crefname{originalequation}{Original equation}{Original equations}
\makeatletter
\newlabel{equation-delta-Z}{{0EIX}{source}{}{sourcepointer-3}{}}
\newlabel{equation-delta-Z@cref}{{[originalequation][0][]0EIX}{[1][0][]source}{}{}{}}
\makeatother
\crefname{originallemma}{original lemma}{original lemmas}
\Crefname{originallemma}{Original lemma}{Original lemmas}
\makeatletter
\newlabel{lemma-discussion}{{0EIY}{source}{}{sourcepointer-4}{}}
\newlabel{lemma-discussion@cref}{{[originallemma][0][]0EIY}{[1][0][]source}{}{}{}}
\makeatother
\crefname{originallemma}{original lemma}{original lemmas}
\Crefname{originallemma}{Original lemma}{Original lemmas}
\makeatletter
\newlabel{coherent-lemma-inverse-systems-pullback}{{0887}{source}{}{sourcepointer-5}{}}
\newlabel{coherent-lemma-inverse-systems-pullback@cref}{{[originallemma][0][]0887}{[1][0][]source}{}{}{}}
\makeatother
\crefname{originallemma}{original lemma}{original lemmas}
\Crefname{originallemma}{Original lemma}{Original lemmas}
\makeatletter
\newlabel{lemma-improvement-formal-coherent-module-better}{{0EJF}{source}{}{sourcepointer-6}{}}
\newlabel{lemma-improvement-formal-coherent-module-better@cref}{{[originallemma][0][]0EJF}{[1][0][]source}{}{}{}}
\makeatother
\crefname{originallemma}{original lemma}{original lemmas}
\Crefname{originallemma}{Original lemma}{Original lemmas}
\makeatletter
\newlabel{divisors-lemma-blowing-up-affine}{{0804}{source}{}{sourcepointer-7}{}}
\newlabel{divisors-lemma-blowing-up-affine@cref}{{[originallemma][0][]0804}{[1][0][]source}{}{}{}}
\makeatother
\crefname{originallemma}{original lemma}{original lemmas}
\Crefname{originallemma}{Original lemma}{Original lemmas}
\makeatletter
\newlabel{more-algebra-lemma-completion-dimension}{{07NV}{source}{}{sourcepointer-8}{}}
\newlabel{more-algebra-lemma-completion-dimension@cref}{{[originallemma][0][]07NV}{[1][0][]source}{}{}{}}
\makeatother
\crefname{originallemma}{original lemma}{original lemmas}
\Crefname{originallemma}{Original lemma}{Original lemmas}
\makeatletter
\newlabel{algebra-lemma-dimension-formula}{{02IJ}{source}{}{sourcepointer-9}{}}
\newlabel{algebra-lemma-dimension-formula@cref}{{[originallemma][0][]02IJ}{[1][0][]source}{}{}{}}
\makeatother
\crefname{originallemma}{original lemma}{original lemmas}
\Crefname{originallemma}{Original lemma}{Original lemmas}
\makeatletter
\newlabel{dualizing-lemma-formal-fibres-gorenstein}{{0BJN}{source}{}{sourcepointer-10}{}}
\newlabel{dualizing-lemma-formal-fibres-gorenstein@cref}{{[originallemma][0][]0BJN}{[1][0][]source}{}{}{}}
\makeatother
\crefname{originallemma}{original lemma}{original lemmas}
\Crefname{originallemma}{Original lemma}{Original lemmas}
\makeatletter
\newlabel{dualizing-lemma-dualizing-gorenstein-formal-fibres}{{0AWY}{source}{}{sourcepointer-11}{}}
\newlabel{dualizing-lemma-dualizing-gorenstein-formal-fibres@cref}{{[originallemma][0][]0AWY}{[1][0][]source}{}{}{}}
\makeatother
\crefname{originalsection}{original section}{original sections}
\Crefname{originalsection}{Original section}{Original sections}
\makeatletter
\newlabel{more-algebra-section-properties-formal-fibres}{{0BIR}{source}{}{sourcepointer-12}{}}
\newlabel{more-algebra-section-properties-formal-fibres@cref}{{[originalsection][0][]0BIR}{[1][0][]source}{}{}{}}
\makeatother
\crefname{originallemma}{original lemma}{original lemmas}
\Crefname{originallemma}{Original lemma}{Original lemmas}
\makeatletter
\newlabel{algebra-lemma-apply-grothendieck-module}{{0338}{source}{}{sourcepointer-13}{}}
\newlabel{algebra-lemma-apply-grothendieck-module@cref}{{[originallemma][0][]0338}{[1][0][]source}{}{}{}}
\makeatother
\crefname{originalproposition}{original proposition}{original propositions}
\Crefname{originalproposition}{Original proposition}{Original propositions}
\makeatletter
\newlabel{more-algebra-proposition-ratliff}{{0AW6}{source}{}{sourcepointer-14}{}}
\newlabel{more-algebra-proposition-ratliff@cref}{{[originalproposition][0][]0AW6}{[1][0][]source}{}{}{}}
\makeatother
\crefname{originallemma}{original lemma}{original lemmas}
\Crefname{originallemma}{Original lemma}{Original lemmas}
\makeatletter
\newlabel{lemma-change-distance-function}{{0EIZ}{source}{}{sourcepointer-15}{}}
\newlabel{lemma-change-distance-function@cref}{{[originallemma][0][]0EIZ}{[1][0][]source}{}{}{}}
\makeatother
\crefname{originallemma}{original lemma}{original lemmas}
\Crefname{originallemma}{Original lemma}{Original lemmas}
\makeatletter
\newlabel{lemma-cd-1-canonical}{{0EJI}{source}{}{sourcepointer-16}{}}
\newlabel{lemma-cd-1-canonical@cref}{{[originallemma][0][]0EJI}{[1][0][]source}{}{}{}}
\makeatother
\crefname{originallemma}{original lemma}{original lemmas}
\Crefname{originallemma}{Original lemma}{Original lemmas}
\makeatletter
\newlabel{lemma-canonically-algebraizable}{{0EIR}{source}{}{sourcepointer-17}{}}
\newlabel{lemma-canonically-algebraizable@cref}{{[originallemma][0][]0EIR}{[1][0][]source}{}{}{}}
\makeatother
\crefname{originallemma}{original lemma}{original lemmas}
\Crefname{originallemma}{Original lemma}{Original lemmas}
\makeatletter
\newlabel{lemma-canonically-extend-base-change}{{0EIS}{source}{}{sourcepointer-18}{}}
\newlabel{lemma-canonically-extend-base-change@cref}{{[originallemma][0][]0EIS}{[1][0][]source}{}{}{}}
\makeatother
\crefname{originallemma}{original lemma}{original lemmas}
\Crefname{originallemma}{Original lemma}{Original lemmas}
\makeatletter
\newlabel{coherent-lemma-existence-projective}{{0885}{source}{}{sourcepointer-19}{}}
\newlabel{coherent-lemma-existence-projective@cref}{{[originallemma][0][]0885}{[1][0][]source}{}{}{}}
\makeatother
\crefname{originallemma}{original lemma}{original lemmas}
\Crefname{originallemma}{Original lemma}{Original lemmas}
\makeatletter
\newlabel{coherent-lemma-existence-easy}{{0889}{source}{}{sourcepointer-20}{}}
\newlabel{coherent-lemma-existence-easy@cref}{{[originallemma][0][]0889}{[1][0][]source}{}{}{}}
\makeatother
\crefname{originallemma}{original lemma}{original lemmas}
\Crefname{originallemma}{Original lemma}{Original lemmas}
\makeatletter
\newlabel{coherent-lemma-inverse-systems-push-pull}{{088B}{source}{}{sourcepointer-21}{}}
\newlabel{coherent-lemma-inverse-systems-push-pull@cref}{{[originallemma][0][]088B}{[1][0][]source}{}{}{}}
\makeatother
\crefname{originallemma}{original lemma}{original lemmas}
\Crefname{originallemma}{Original lemma}{Original lemmas}
\makeatletter
\newlabel{lemma-map-kernel-cokernel-on-closed}{{0EIU}{source}{}{sourcepointer-22}{}}
\newlabel{lemma-map-kernel-cokernel-on-closed@cref}{{[originallemma][0][]0EIU}{[1][0][]source}{}{}{}}
\makeatother
\crefname{originallemma}{original lemma}{original lemmas}
\Crefname{originallemma}{Original lemma}{Original lemmas}
\makeatletter
\newlabel{lemma-algebraizable}{{0EIN}{source}{}{sourcepointer-23}{}}
\newlabel{lemma-algebraizable@cref}{{[originallemma][0][]0EIN}{[1][0][]source}{}{}{}}
\makeatother
\crefname{originallemma}{original lemma}{original lemmas}
\Crefname{originallemma}{Original lemma}{Original lemmas}
\makeatletter
\newlabel{dualizing-lemma-ubiquity-dualizing}{{0BFR}{source}{}{sourcepointer-24}{}}
\newlabel{dualizing-lemma-ubiquity-dualizing@cref}{{[originallemma][0][]0BFR}{[1][0][]source}{}{}{}}
\makeatother
\crefname{originallemma}{original lemma}{original lemmas}
\Crefname{originallemma}{Original lemma}{Original lemmas}
\makeatletter
\newlabel{divisors-lemma-flat-base-change-blowing-up}{{0805}{source}{}{sourcepointer-25}{}}
\newlabel{divisors-lemma-flat-base-change-blowing-up@cref}{{[originallemma][0][]0805}{[1][0][]source}{}{}{}}
\makeatother
\title{Lemma blowup in algebraization.tex}
\author{The Stacks Project Authors}
\date{stacks-project:a04446e57ec1fbc252a871afcec7752fb2807b14; GNU Free Documentation License 1.2 (as published in the Stacks Project COPYING)}
\begin{document}
\maketitle
\noindent\textit{Editorial note.} All mathematical statement and proof text is reproduced verbatim from the author's chapter algebraization.tex of the Stacks Project, fixed at commit a04446e57ec1fbc252a871afcec7752fb2807b14. The pinned preamble is supplied as its own file; the wrapper class, the xr-hyper wiring and the \textbackslash{}externaldocument list are omitted because the assembler supplies the class.\par\medskip
\noindent\textit{Editorial note.} Numbering inside this standalone excerpt is editorial, not the full Stacks book print numbering.\par\medskip
\noindent\textit{Editorial note.} The candidate is from the corrected band-I follow-on screen (line C) of the extractor's stacks-candidates line: its author label is absent from the live master in every form searched and it belongs to none of the earlier judged pools. The H2 judgement and its reason were made in this module by reading the author's own statement and complete proof in the fixed checkout. Original candidate key: stacks:algebraization.tex:lemma-blowup.\par\medskip
\noindent\textit{Editorial note.} Reference labels in the proof resolve to navigation pointers to the original author source; these pointers do not reproduce proofs of the prerequisites.\par\medskip
\noindent\textit{Allowed standard prerequisites (editorial).}\begin{itemize}
\item Prerequisites named by the author inside the counted span. The counted span is reproduced verbatim, including the author's own citations. Its prerequisites are the earlier results of the same fixed Stacks Project version that the author cites there and in the chapter's lead-in.
\end{itemize}\medskip
\ifdefined\hypertarget\hypertarget{sourcepointer-1}{}\fi
\noindent\textit{Source reference.} Original reference 0EHC: sources/algebraization.tex, source lines 4983--4983. Navigation pointer used in the counted span; the span contains the pointer only and the cited result is not reproduced inside it. Registered as a non-load-bearing navigation stub in the established form for Stacks material. No bibliography citation is involved.\par\medskip
\ifdefined\hypertarget\hypertarget{sourcepointer-2}{}\fi
\noindent\textit{Source reference.} Original definition 0EJ3: sources/algebraization.tex, source lines 5928--5928. Navigation pointer used in the counted span; the span contains the pointer only and the cited result is not reproduced inside it. Registered as a non-load-bearing navigation stub in the established form for Stacks material. No bibliography citation is involved.\par\medskip
\ifdefined\hypertarget\hypertarget{sourcepointer-3}{}\fi
\noindent\textit{Source reference.} Original equation 0EIX: sources/algebraization.tex, source lines 5589--5589. Navigation pointer used in the counted span; the span contains the pointer only and the cited result is not reproduced inside it. Registered as a non-load-bearing navigation stub in the established form for Stacks material. No bibliography citation is involved.\par\medskip
\ifdefined\hypertarget\hypertarget{sourcepointer-4}{}\fi
\noindent\textit{Source reference.} Original lemma 0EIY: sources/algebraization.tex, source lines 5622--5622. Navigation pointer used in the counted span; the span contains the pointer only and the cited result is not reproduced inside it. Registered as a non-load-bearing navigation stub in the established form for Stacks material. No bibliography citation is involved.\par\medskip
\ifdefined\hypertarget\hypertarget{sourcepointer-5}{}\fi
\noindent\textit{Source reference.} Original lemma 0887: sources/coherent.tex, source lines 6631--6631. Navigation pointer used in the counted span; the span contains the pointer only and the cited result is not reproduced inside it. Registered as a non-load-bearing navigation stub in the established form for Stacks material. No bibliography citation is involved.\par\medskip
\ifdefined\hypertarget\hypertarget{sourcepointer-6}{}\fi
\noindent\textit{Source reference.} Original lemma 0EJF: sources/algebraization.tex, source lines 6952--6952. Navigation pointer used in the counted span; the span contains the pointer only and the cited result is not reproduced inside it. Registered as a non-load-bearing navigation stub in the established form for Stacks material. No bibliography citation is involved.\par\medskip
\ifdefined\hypertarget\hypertarget{sourcepointer-7}{}\fi
\noindent\textit{Source reference.} Original lemma 0804: sources/divisors.tex, source lines 8322--8322. Navigation pointer used in the counted span; the span contains the pointer only and the cited result is not reproduced inside it. Registered as a non-load-bearing navigation stub in the established form for Stacks material. No bibliography citation is involved.\par\medskip
\ifdefined\hypertarget\hypertarget{sourcepointer-8}{}\fi
\noindent\textit{Source reference.} Original lemma 07NV: sources/more-algebra.tex, source lines 10624--10624. Navigation pointer used in the counted span; the span contains the pointer only and the cited result is not reproduced inside it. Registered as a non-load-bearing navigation stub in the established form for Stacks material. No bibliography citation is involved.\par\medskip
\ifdefined\hypertarget\hypertarget{sourcepointer-9}{}\fi
\noindent\textit{Source reference.} Original lemma 02IJ: sources/algebra.tex, source lines 27520--27520. Navigation pointer used in the counted span; the span contains the pointer only and the cited result is not reproduced inside it. Registered as a non-load-bearing navigation stub in the established form for Stacks material. No bibliography citation is involved.\par\medskip
\ifdefined\hypertarget\hypertarget{sourcepointer-10}{}\fi
\noindent\textit{Source reference.} Original lemma 0BJN: sources/dualizing.tex, source lines 4191--4191. Navigation pointer used in the counted span; the span contains the pointer only and the cited result is not reproduced inside it. Registered as a non-load-bearing navigation stub in the established form for Stacks material. No bibliography citation is involved.\par\medskip
\ifdefined\hypertarget\hypertarget{sourcepointer-11}{}\fi
\noindent\textit{Source reference.} Original lemma 0AWY: sources/dualizing.tex, source lines 4235--4235. Navigation pointer used in the counted span; the span contains the pointer only and the cited result is not reproduced inside it. Registered as a non-load-bearing navigation stub in the established form for Stacks material. No bibliography citation is involved.\par\medskip
\ifdefined\hypertarget\hypertarget{sourcepointer-12}{}\fi
\noindent\textit{Source reference.} Original section 0BIR: sources/more-algebra.tex, source lines 12814--12814. Navigation pointer used in the counted span; the span contains the pointer only and the cited result is not reproduced inside it. Registered as a non-load-bearing navigation stub in the established form for Stacks material. No bibliography citation is involved.\par\medskip
\ifdefined\hypertarget\hypertarget{sourcepointer-13}{}\fi
\noindent\textit{Source reference.} Original lemma 0338: sources/algebra.tex, source lines 46303--46303. Navigation pointer used in the counted span; the span contains the pointer only and the cited result is not reproduced inside it. Registered as a non-load-bearing navigation stub in the established form for Stacks material. No bibliography citation is involved.\par\medskip
\ifdefined\hypertarget\hypertarget{sourcepointer-14}{}\fi
\noindent\textit{Source reference.} Original proposition 0AW6: sources/more-algebra.tex, source lines 32484--32484. Navigation pointer used in the counted span; the span contains the pointer only and the cited result is not reproduced inside it. Registered as a non-load-bearing navigation stub in the established form for Stacks material. No bibliography citation is involved.\par\medskip
\ifdefined\hypertarget\hypertarget{sourcepointer-15}{}\fi
\noindent\textit{Source reference.} Original lemma 0EIZ: sources/algebraization.tex, source lines 5760--5760. Navigation pointer used in the counted span; the span contains the pointer only and the cited result is not reproduced inside it. Registered as a non-load-bearing navigation stub in the established form for Stacks material. No bibliography citation is involved.\par\medskip
\ifdefined\hypertarget\hypertarget{sourcepointer-16}{}\fi
\noindent\textit{Source reference.} Original lemma 0EJI: sources/algebraization.tex, source lines 7129--7129. Navigation pointer used in the counted span; the span contains the pointer only and the cited result is not reproduced inside it. Registered as a non-load-bearing navigation stub in the established form for Stacks material. No bibliography citation is involved.\par\medskip
\ifdefined\hypertarget\hypertarget{sourcepointer-17}{}\fi
\noindent\textit{Source reference.} Original lemma 0EIR: sources/algebraization.tex, source lines 5249--5249. Navigation pointer used in the counted span; the span contains the pointer only and the cited result is not reproduced inside it. Registered as a non-load-bearing navigation stub in the established form for Stacks material. No bibliography citation is involved.\par\medskip
\ifdefined\hypertarget\hypertarget{sourcepointer-18}{}\fi
\noindent\textit{Source reference.} Original lemma 0EIS: sources/algebraization.tex, source lines 5297--5297. Navigation pointer used in the counted span; the span contains the pointer only and the cited result is not reproduced inside it. Registered as a non-load-bearing navigation stub in the established form for Stacks material. No bibliography citation is involved.\par\medskip
\ifdefined\hypertarget\hypertarget{sourcepointer-19}{}\fi
\noindent\textit{Source reference.} Original lemma 0885: sources/coherent.tex, source lines 6812--6812. Navigation pointer used in the counted span; the span contains the pointer only and the cited result is not reproduced inside it. Registered as a non-load-bearing navigation stub in the established form for Stacks material. No bibliography citation is involved.\par\medskip
\ifdefined\hypertarget\hypertarget{sourcepointer-20}{}\fi
\noindent\textit{Source reference.} Original lemma 0889: sources/coherent.tex, source lines 6487--6487. Navigation pointer used in the counted span; the span contains the pointer only and the cited result is not reproduced inside it. Registered as a non-load-bearing navigation stub in the established form for Stacks material. No bibliography citation is involved.\par\medskip
\ifdefined\hypertarget\hypertarget{sourcepointer-21}{}\fi
\noindent\textit{Source reference.} Original lemma 088B: sources/coherent.tex, source lines 7025--7025. Navigation pointer used in the counted span; the span contains the pointer only and the cited result is not reproduced inside it. Registered as a non-load-bearing navigation stub in the established form for Stacks material. No bibliography citation is involved.\par\medskip
\ifdefined\hypertarget\hypertarget{sourcepointer-22}{}\fi
\noindent\textit{Source reference.} Original lemma 0EIU: sources/algebraization.tex, source lines 5458--5458. Navigation pointer used in the counted span; the span contains the pointer only and the cited result is not reproduced inside it. Registered as a non-load-bearing navigation stub in the established form for Stacks material. No bibliography citation is involved.\par\medskip
\ifdefined\hypertarget\hypertarget{sourcepointer-23}{}\fi
\noindent\textit{Source reference.} Original lemma 0EIN: sources/algebraization.tex, source lines 5171--5171. Navigation pointer used in the counted span; the span contains the pointer only and the cited result is not reproduced inside it. Registered as a non-load-bearing navigation stub in the established form for Stacks material. No bibliography citation is involved.\par\medskip
\ifdefined\hypertarget\hypertarget{sourcepointer-24}{}\fi
\noindent\textit{Source reference.} Original lemma 0BFR: sources/dualizing.tex, source lines 4126--4126. Navigation pointer used in the counted span; the span contains the pointer only and the cited result is not reproduced inside it. Registered as a non-load-bearing navigation stub in the established form for Stacks material. No bibliography citation is involved.\par\medskip
\ifdefined\hypertarget\hypertarget{sourcepointer-25}{}\fi
\noindent\textit{Source reference.} Original lemma 0805: sources/divisors.tex, source lines 8351--8351. Navigation pointer used in the counted span; the span contains the pointer only and the cited result is not reproduced inside it. Registered as a non-load-bearing navigation stub in the established form for Stacks material. No bibliography citation is involved.\par\medskip
\begin{lemma}
\label{lemma-blowup}
In Situation \ref{situation-algebraize} let $(\mathcal{F}_n)$
be an object of $\textit{Coh}(U, I\mathcal{O}_U)$. Assume
\begin{enumerate}
\item $A$ has a dualizing complex,
\item all fibres of the blowing up $b : X' \to X$ of $I$
have dimension $\leq d - 1$,
\item one of the following is true
\begin{enumerate}
\item $(\mathcal{F}_n)$ satisfies the $(d + 1, d + 2)$-inequalities
(Definition \ref{definition-s-d-inequalities}), or
\item for $y \in U \cap Y$ and a prime
$\mathfrak p \subset \mathcal{O}_{X, y}^\wedge$ with
$\mathfrak p \not \in V(I\mathcal{O}_{X, y}^\wedge)$
we have
$$
\text{depth}((\mathcal{F}^\wedge_y)_\mathfrak p) +
\dim(\mathcal{O}_{X, y}^\wedge/\mathfrak p) + \delta^Y_Z(y) > d + 2
$$
\end{enumerate}
\end{enumerate}
Then $(\mathcal{F}_n)$ extends to $X$.
\end{lemma}

\begin{proof}
Let $Y' \subset X'$ be the exceptional divisor.
Let $Z' \subset Y'$ be the inverse image of $Z \subset Y$.
Then $U' = X' \setminus Z'$ is the inverse image of $U$.
With $\delta^{Y'}_{Z'}$ as in (\ref{equation-delta-Z}) we set
$$
T' = \{y' \in Y' \mid \delta^{Y'}_{Z'}(y') = 1\}
\subset
T = \{y' \in Y' \mid \delta^{Y'}_{Z'}(y') = 1\text{ or }2\}
$$
By Lemma \ref{lemma-discussion} parts (1) and (2) these are specialization
stable subsets of $U' \cap Y' = Y' \setminus Z'$. Consider the
object $(b|_{U'}^*\mathcal{F}_n)$ of $\textit{Coh}(U', I\mathcal{O}_{U'})$,
see Cohomology of Schemes, Lemma \ref{coherent-lemma-inverse-systems-pullback}.
For $y' \in U' \cap Y'$ let us denote
$$
\mathcal{F}_{y'}^\wedge = \lim (b|_{U'}^*\mathcal{F}_n)_{y'}
$$
the ``stalk'' of this pullback at $y'$. We claim that conditions
(a), (b), (c), (d), and (e) of
Lemma \ref{lemma-improvement-formal-coherent-module-better}
hold for the object $(b|_{U'}^*\mathcal{F}_n)$ on $U'$ with $d$
replaced by $1$ and the subsets $T' \subset T \subset U' \cap Y'$.
Condition (a) holds because $Y'$ is an effective Cartier divisor
and hence locally cut out by $1$ equation. Condition (e) holds
by Lemma \ref{lemma-discussion} part (7).
To prove (b), (c), and (d) we need some preparation.

\medskip\noindent
Let $y' \in U' \cap Y'$ and let
$\mathfrak p' \subset \mathcal{O}_{X', y'}^\wedge$
be a prime ideal not contained in $V(I\mathcal{O}_{X', y'}^\wedge)$.
Denote $y = b(y') \in U \cap Y$. Choose $f \in I$ such that
$y'$ is contained in the spectrum of the affine blowup algebra
$A[\frac{I}{f}]$, see Divisors, Lemma \ref{divisors-lemma-blowing-up-affine}.
For any $A$-algebra $B$ denote $B' = B[\frac{IB}{f}]$ the corresponding affine
blowup algebra. Denote $I$-adic completion by ${\ }^\wedge$.
By our choice of $f$ we get a ring map
$(\mathcal{O}_{X, y}^\wedge)' \to \mathcal{O}_{X', y'}^\wedge$.
If we let $\mathfrak q' \subset (\mathcal{O}_{X, y}^\wedge)'$
be the inverse image of $\mathfrak m_{y'}^\wedge$, then
we see that
$((\mathcal{O}_{X, y}^\wedge)'_{\mathfrak q'})^\wedge =
\mathcal{O}_{X', y'}^\wedge$.
Let $\mathfrak p \subset \mathcal{O}_{X, y}^\wedge$ be the corresponding
prime. At this point we have a commutative diagram
$$
\xymatrix{
\mathcal{O}_{X, y}^\wedge \ar[d] \ar[r] &
(\mathcal{O}_{X, y}^\wedge)' \ar[d]_\alpha \ar[r] &
(\mathcal{O}_{X, y}^\wedge)'_{\mathfrak q'} \ar[d] \ar[r]_\beta &
\mathcal{O}_{X', y'}^\wedge \ar[d] \\
\mathcal{O}_{X, y}^\wedge/\mathfrak p \ar[r] &
(\mathcal{O}_{X, y}^\wedge/\mathfrak p)' \ar[r] &
(\mathcal{O}_{X, y}^\wedge/\mathfrak p)'_{\mathfrak q'} \ar[r]^\gamma &
((\mathcal{O}_{X, y}^\wedge/\mathfrak p)'_{\mathfrak q'})^\wedge \ar[d] \\
 & & &
\mathcal{O}_{X', y'}^\wedge/\mathfrak p'
}
$$
whose vertical arrows are surjective. By
More on Algebra, Lemma \ref{more-algebra-lemma-completion-dimension}
and the dimension formula
(Algebra, Lemma \ref{algebra-lemma-dimension-formula})
we have
$$
\dim(((\mathcal{O}_{X, y}^\wedge/\mathfrak p)'_{\mathfrak q'})^\wedge) =
\dim((\mathcal{O}_{X, y}^\wedge/\mathfrak p)'_{\mathfrak q'}) =
\dim(\mathcal{O}_{X, y}^\wedge/\mathfrak p)
- \text{trdeg}(\kappa(y')/\kappa(y))
$$
Tracing through the definitions of pullbacks, stalks, localizations,
and completions we find
$$
(\mathcal{F}_y^\wedge)_{\mathfrak p}
\otimes_{(\mathcal{O}_{X, y}^\wedge)_\mathfrak p}
(\mathcal{O}_{X', y'}^\wedge)_{\mathfrak p'}
=
(\mathcal{F}_{y'}^\wedge)_{\mathfrak p'}
$$
Details omitted. The ring maps $\beta$ and $\gamma$ in the diagram
are flat with Gorenstein (hence Cohen-Macaulay) fibres, as these are
completions of rings having a dualizing complex. See
Dualizing Complexes, Lemmas
\ref{dualizing-lemma-formal-fibres-gorenstein} and
\ref{dualizing-lemma-dualizing-gorenstein-formal-fibres}
and the discussion in More on Algebra, Section
\ref{more-algebra-section-properties-formal-fibres}.
Observe that $(\mathcal{O}_{X, y}^\wedge)_\mathfrak p =
(\mathcal{O}_{X, y}^\wedge)'_{\tilde{\mathfrak p}}$
where $\tilde{\mathfrak p}$ is the kernel of $\alpha$
in the diagram. On the other hand,
$(\mathcal{O}_{X, y}^\wedge)'_{\tilde{\mathfrak p}}
\to (\mathcal{O}_{X', y'}^\wedge)_{\mathfrak p'}$
is flat with CM fibres by the above. Whence
$(\mathcal{O}_{X, y}^\wedge)_\mathfrak p  \to
(\mathcal{O}_{X', y'}^\wedge)_{\mathfrak p'}$ is flat with CM fibres.
Using Algebra, Lemma \ref{algebra-lemma-apply-grothendieck-module}
we see that
$$
\text{depth}((\mathcal{F}_{y'}^\wedge)_{\mathfrak p'}) =
\text{depth}((\mathcal{F}_y^\wedge)_{\mathfrak p}) +
\dim(F_\mathfrak r)
$$
where $F$ is the generic formal fibre of
$(\mathcal{O}_{X, y}^\wedge/\mathfrak p)'_{\mathfrak q'}$
and $\mathfrak r$ is the prime corresponding to $\mathfrak p'$.
Since $(\mathcal{O}_{X, y}^\wedge/\mathfrak p)'_{\mathfrak q'}$
is a universally catenary local domain, its $I$-adic completion
is equidimensional and (universally) catenary by Ratliff's theorem
(More on Algebra, Proposition \ref{more-algebra-proposition-ratliff}).
It then follows that
$$
\dim(((\mathcal{O}_{X, y}^\wedge/\mathfrak p)'_{\mathfrak q'})^\wedge) =
\dim(F_\mathfrak r) + \dim(\mathcal{O}_{X', y'}^\wedge/\mathfrak p')
$$
Combined with Lemma \ref{lemma-change-distance-function}
we get
\begin{equation}
\label{equation-one}
\begin{aligned}
&
\text{depth}((\mathcal{F}_{y'}^\wedge)_{\mathfrak p'}) +
\delta^{Y'}_{Z'}(y') \\
& =
\text{depth}((\mathcal{F}_y^\wedge)_{\mathfrak p}) +
\dim(F_\mathfrak r) + \delta^{Y'}_{Z'}(y') \\
& \geq
\text{depth}((\mathcal{F}_y^\wedge)_{\mathfrak p}) + \delta^Y_Z(y) +
\dim(F_\mathfrak r) + \text{trdeg}(\kappa(y')/\kappa(y)) - (d - 1) \\
& =
\text{depth}((\mathcal{F}_y^\wedge)_{\mathfrak p}) + \delta^Y_Z(y) - (d - 1)
+ \dim(\mathcal{O}_{X, y}^\wedge/\mathfrak p) -
\dim(\mathcal{O}_{X', y'}^\wedge/\mathfrak p')
\end{aligned}
\end{equation}
Please keep in mind that
$\dim(\mathcal{O}_{X, y}^\wedge/\mathfrak p) \geq
\dim(\mathcal{O}_{X', y'}^\wedge/\mathfrak p')$. Rewriting this we get
\begin{equation}
\label{equation-two}
\begin{aligned}
&
\text{depth}((\mathcal{F}_{y'}^\wedge)_{\mathfrak p'}) +
\dim(\mathcal{O}_{X', y'}^\wedge/\mathfrak p') +
\delta^{Y'}_{Z'}(y') \\
& \geq
\text{depth}((\mathcal{F}_y^\wedge)_{\mathfrak p}) +
\dim(\mathcal{O}_{X, y}^\wedge/\mathfrak p) +
\delta^Y_Z(y) - (d - 1)
\end{aligned}
\end{equation}
This inequality will allow us to check the remaining conditions.

\medskip\noindent
Conditions (b) and (d) of
Lemma \ref{lemma-improvement-formal-coherent-module-better}. Assume
$V(\mathfrak p') \cap V(I\mathcal{O}_{X', y'}^\wedge) =
\{\mathfrak m_{y'}^\wedge\}$.
This implies that $\dim(\mathcal{O}_{X', y'}^\wedge/\mathfrak p') = 1$
because $Z'$ is an effective Cartier divisor.
The combination of (b) and (d) is equivalent with
$$
\text{depth}((\mathcal{F}_{y'}^\wedge)_{\mathfrak p'}) + \delta^{Y'}_{Z'}(y')
> 2
$$
If $(\mathcal{F}_n)$ satisfies the inequalities in (3)(b)
then we immediately conclude this is true by applying (\ref{equation-two}).
If $(\mathcal{F}_n)$ satisfies (3)(a), i.e., the
$(d + 1, d + 2)$-inequalities, then we see that in any case
$$
\text{depth}((\mathcal{F}_y^\wedge)_{\mathfrak p}) + \delta^Y_Z(y)
\geq d + 1
\quad\text{or}\quad
\text{depth}((\mathcal{F}_y^\wedge)_{\mathfrak p}) +
\dim(\mathcal{O}_{X, y}^\wedge/\mathfrak p) +
\delta^Y_Z(y) > d + 2
$$
Looking at (\ref{equation-one}) and (\ref{equation-two}) above this gives
what we want except possibly if
$\dim(\mathcal{O}_{X, y}^\wedge/\mathfrak p) = 1$.
However, if $\dim(\mathcal{O}_{X, y}^\wedge/\mathfrak p) = 1$, then we have
$V(\mathfrak p) \cap V(I\mathcal{O}_{X, y}^\wedge) = \{\mathfrak m_y^\wedge\}$
and we see that actually
$$
\text{depth}((\mathcal{F}_y^\wedge)_{\mathfrak p}) + \delta^Y_Z(y) > d + 1
$$
as $(\mathcal{F}_n)$ satisfies the $(d + 1, d + 2)$-inequalities and we
conclude again.

\medskip\noindent
Condition (c) of
Lemma \ref{lemma-improvement-formal-coherent-module-better}. Assume
$V(\mathfrak p') \cap V(I\mathcal{O}_{X', y'}^\wedge) \not =
\{\mathfrak m_{y'}^\wedge\}$. Then condition (c) is equivalent to
$$
\text{depth}((\mathcal{F}_{y'}^\wedge)_{\mathfrak p'}) + \delta^{Y'}_{Z'}(y')
\geq 2
\quad\text{or}\quad
\text{depth}((\mathcal{F}_{y'}^\wedge)_{\mathfrak p'}) +
\dim(\mathcal{O}_{X', y'}^\wedge/\mathfrak p') +
\delta^{Y'}_{Z'}(y') > 3
$$
If $(\mathcal{F}_n)$ satisfies the inequalities in (3)(b)
then we see the second of the two displayed inequalities holds true
by applying (\ref{equation-two}). If $(\mathcal{F}_n)$ satisfies (3)(a), i.e.,
the $(d + 1, d + 2)$-inequalities, then this follows immediately from
(\ref{equation-one}) and (\ref{equation-two}).
This finishes the proof of our claim.

\medskip\noindent
Choose $(b|_{U'}^*\mathcal{F}_n) \to (\mathcal{F}_n'')$
and $(\mathcal{H}_n)$ in $\textit{Coh}(U', I\mathcal{O}_{U'})$
as in Lemma \ref{lemma-improvement-formal-coherent-module-better}.
For any affine open $W \subset X'$ observe that
$\delta^{W \cap Y'}_{W \cap Z'}(y') \geq \delta^{Y'}_{Z'}(y')$ by
Lemma \ref{lemma-discussion} part (6). Hence we see that
$(\mathcal{H}_n|_W)$ satisfies the assumptions of
Lemma \ref{lemma-cd-1-canonical}.
Thus $(\mathcal{H}_n|_W)$ extends canonically to $W$.
Let $(\mathcal{G}_{W, n})$ in $\textit{Coh}(W, I\mathcal{O}_W)$
be the canonical extension as in
Lemma \ref{lemma-canonically-algebraizable}.
By Lemma \ref{lemma-canonically-extend-base-change}
we see that for $W' \subset W$ there is a unique isomorphism
$$
(\mathcal{G}_{W, n}|_{W'}) \longrightarrow
(\mathcal{G}_{W', n})
$$
compatible with the given isomorphisms
$(\mathcal{G}_{W, n}|_{W \cap U}) \cong (\mathcal{H}_n|_{W \cap U})$.
We conclude that there exists an object
$(\mathcal{G}_n)$ of $\textit{Coh}(X', I\mathcal{O}_{X'})$
whose restriction to $U$ is isomorphic to $(\mathcal{H}_n)$.

\medskip\noindent
If $A$ is $I$-adically complete we can finish the proof as follows.
By Grothendieck's existence theorem
(Cohomology of Schemes, Lemma \ref{coherent-lemma-existence-projective})
we see that $(\mathcal{G}_n)$ is the completion of a coherent
$\mathcal{O}_{X'}$-module. Then by
Cohomology of Schemes, Lemma \ref{coherent-lemma-existence-easy}
we see that $(b|_{U'}^*\mathcal{F}_n)$
is the completion of a coherent $\mathcal{O}_{U'}$-module
$\mathcal{F}'$.
By Cohomology of Schemes, Lemma \ref{coherent-lemma-inverse-systems-push-pull}
we see that there is a map
$$
(\mathcal{F}_n) \longrightarrow ((b|_{U'})_*\mathcal{F}')^\wedge
$$
whose kernel and cokernel is annihilated by a power of $I$.
Then finally, we win by applying
Lemma \ref{lemma-map-kernel-cokernel-on-closed}.

\medskip\noindent
If $A$ is not complete, then, before starting the proof, we may replace $A$
by its completion, see Lemma \ref{lemma-algebraizable}.
After completion the assumptions still hold: this is immediate
for condition (3), follows from
Dualizing Complexes, Lemma \ref{dualizing-lemma-ubiquity-dualizing}
for condition (1), and from
Divisors, Lemma \ref{divisors-lemma-flat-base-change-blowing-up}
for condition (2).
Thus the complete case implies the general case.
\end{proof}

\end{document}
